% ===================================================================
\section{Glossary of terms}\label{sec:glossary}
% ===================================================================
{\small
\setlength{\columnsep}{1.4em}
\begin{multicols}{2}
\begin{description}[leftmargin=0pt,itemsep=2pt,font=\bfseries]
\item[Adaptation] A trait shaped by natural selection that suits the environment; Darwin's Inf 3.
\item[Adaptive force (Lamarck)] Force fitting organisms to local conditions by use/disuse and inheritance of acquired characters.
\item[Age pyramid] Bar chart of age-class shares, males and females back to back; base broadens with $\lambda$.
\item[Allee effect] Per-capita growth falling as density falls; strong form has an extinction threshold $A$.
\item[Allele frequency] Share of gene copies of an allele: $p=P+Q/2$ (always true).
\item[Alternative stable states] Two or more stable states under the same conditions (clear vs turbid shallow lake).
\item[Antagonistic pleiotropy] Ageing theory: alleles with early benefits and late costs are favoured.
\item[Aposematism] Conspicuous warning signal of a defended prey, letting predators learn.
\item[Batesian mimicry] Harmless mimic copies a defended model's signal; protected only while rarer than the model.
\item[Bifurcation diagram] Long-run values visited against a parameter; number of points = period.
\item[Biomanipulation] Reducing planktivorous fish so grazers clear the water; a push across a threshold.
\item[Blending inheritance] Offspring = mid-parent value; halves heritable variance each generation (Jenkin).
\item[Buffer zone] Partly used land around a reserve core; holds prey, dispersal routes, distance from people.
\item[Carrying capacity ($K$)] Density at which the realised per-capita growth rate is zero.
\item[Catalepsis] Playing dead (thanatosis), e.g.\ the opossum.
\item[Causal loop diagram] Network of signed effects; the loop sign is the product of link signs.
\item[Character displacement] Related species differ \emph{more} where they co-occur than where apart.
\item[Chaos (deterministic)] Bounded, aperiodic dynamics with sensitive dependence on initial conditions; no noise involved.
\item[Clinograde / orthograde] Summer O$_2$ profile falling to zero with depth (eutrophic) / staying high (oligotrophic).
\item[Cohort] Individuals born in the same period, followed together.
\item[Cohort life table] Record of a cohort's survivorship $l_x$ and fecundity $m_x$ by age.
\item[Competition] Use or defence of a resource that reduces its availability to others; non-trophic.
\item[Competitive exclusion principle] No two species with the same niche (one limiting resource) can coexist (Gause).
\item[Complexifying force (Lamarck)] Force moving lineages up a master plan toward greater complexity.
\item[Contest / scramble competition] Resource shared out so more parents never mean fewer offspring / split so finely that crowding lowers output (overcompensation).
\item[Crypsis] Avoiding detection by matching the background.
\item[Cultural eutrophication] Nutrient enrichment of waters by human activity.
\item[Demographic stochasticity] Chance variation in births, deaths, sex ratio; relative size $\sim1/\sqrt N$.
\item[Density dependence] Per-capita rates changing with density (negative: crowding; positive: Allee).
\item[Descent with modification] All species are branches of one tree, changed from common ancestors.
\item[Detection probability ($p$)] Chance an animal present is counted; $\hat N=C/\hat p$.
\item[Directional / stabilizing / disruptive selection] Favours one extreme (mean shifts) / intermediates (variance falls) / both extremes (variance rises).
\item[Dominant eigenvalue ($\lambda$)] Largest-modulus eigenvalue of a projection matrix; long-run growth factor.
\item[Dystrophic lake] Low conductivity, brown humic water.
\item[Ecology] Study of interactions of organisms with their biotic and abiotic environment, explaining distribution and abundance.
\item[Effective population size ($N_e$)] Size of an ideal population with the same drift rate; usually $<$ census $N$.
\item[Elasticity] Proportional change in $\lambda$ per proportional change in a matrix entry; elasticities sum to 1.
\item[Emergence] Properties of a whole not found in its parts (a migration, a traffic jam).
\item[Environmental stochasticity] Shocks affecting all individuals together; does not average out with $N$.
\item[Epigenetic inheritance] Transmission of gene-expression states (e.g.\ methylation) without DNA sequence change.
\item[Ergodicity (demographic)] A primitive matrix population forgets its initial structure.
\item[Eutrophic] Highly productive, green, often with anoxic hypolimnion; $P_{tot}$ 30--100\,µg\,L$^{-1}$.
\item[Extended Evolutionary Synthesis] Proposed extension of the MS: epigenetics, inheritance of some acquired states, HGT, evo-devo.
\item[Fitness] Relative contribution to the next generation (survival $\times$ reproduction) of a genotype or phenotype.
\item[Fractal] Self-similar at many scales (the bifurcation diagram).
\item[Functional response] Prey eaten per predator as a function of prey density (Holling II saturates at $1/h$).
\item[Fundamental / realised niche] Where a species could persist alone / where it actually does with enemies and competitors.
\item[Genetic drift] Random change in allele frequency from sampling; strongest in small populations.
\item[Genetic stochasticity] Drift, loss of heterozygosity and inbreeding depression in small populations.
\item[Hardy--Weinberg equilibrium] $p^2:2pq:q^2$ with constant frequencies under no evolutionary forces.
\item[Heritability ($h^2$)] Fraction of trait variation that is (additively) genetic; sets response $R=h^2S$.
\item[Horizontal gene transfer] Genes moving between contemporaneous organisms, even across species.
\item[Hypolimnion / epilimnion] Cold dark bottom layer / warm mixed top layer of a stratified lake.
\item[Hysteresis] The return path differs from the outward path; recovery needs a larger reversal.
\item[Inbreeding depression] Fitness loss from mating among relatives exposing deleterious recessives.
\item[Internal loading] Phosphorus release from (anoxic) sediments; makes eutrophication self-accelerating.
\item[Intrinsic / extrinsic cause of cycles] Generated by interacting populations / imposed from outside (weather, sunspots).
\item[Irreducible / primitive matrix] Life-cycle graph connected / loop lengths with gcd 1; together they guarantee a SAD.
\item[Isocline (zero-growth)] Curve where one species' growth rate is zero.
\item[Lattice effect] Integer counts with a bound give finitely many states, so chaos becomes cycles or extinction.
\item[Lefkovitch matrix] Stage-structured projection matrix with stay ($P$), grow ($G$) and fecundity ($F$) entries.
\item[Leslie matrix] Age-structured projection matrix: fecundities in row 1, survivals on the sub-diagonal.
\item[Limit cycle] Isolated periodic orbit with fixed amplitude; attracts nearby orbits (unlike neutral L--V cycles).
\item[Lyapunov exponent] Average exponential rate of divergence of nearby orbits; $>0$ means chaos.
\item[Macroevolution / microevolution] Change above the species level over long times / change in trait values or allele frequencies within populations.
\item[Mark--recapture] Estimating $N$ from recapture rates of marked animals ($\hat N=Mn/m$).
\item[Metapopulation] Set of local populations linked by dispersal, persisting through extinction and recolonisation.
\item[Minimum dynamic area] Smallest reserve area holding a mosaic of disturbance and succession patches.
\item[Minimum viable population (MVP)] Smallest size giving an acceptable persistence probability over a stated time.
\item[Modelling relation (Rosen)] Encode nature into a formal system, infer, decode; good if it matches causal entailment.
\item[Modern Synthesis] 1930s--50s union of Darwin and Mendelian genetics: selection on particulate variation; micro $\to$ macro.
\item[Müllerian mimicry] Several defended species share one warning signal.
\item[Mutation accumulation] Ageing theory: late-acting deleterious mutations escape selection and accumulate.
\item[Natural selection] Differential survival and reproduction of individuals due to phenotypic differences.
\item[Net reproductive rate ($R_0$)] Mean lifetime offspring per newborn, $\sum l_xm_x$.
\item[Neutral stability] Displaced system stays at its new state (HW allele frequency; L--V orbits).
\item[Niche (persistence / utilisation)] Conditions allowing positive growth (Hutchinson) / what the species actually uses (MacArthur, Levins).
\item[Niche partitioning] Separation of resource use in space, time, morphology or conditions, allowing coexistence.
\item[Oligotrophic] Low productivity, clear water, oxygenated depths; $P_{tot}$ 5--10\,µg\,L$^{-1}$.
\item[Overcompensation] Crowding cuts output so much that the population overshoots its equilibrium.
\item[Overdominance / underdominance] Heterozygote fitter (stable polymorphism) / less fit (unstable equilibrium).
\item[Period doubling] Successive births of 2-, 4-, 8-cycles as a parameter rises; route to chaos.
\item[Phase plane] Plot of one population against another, with isoclines and trajectories.
\item[Population] Similar individuals inhabiting a defined space at a given time.
\item[Population thinking] Every individual is unique; the mean is an abstraction (vs essentialism).
\item[Preadaptation] Existing trait co-opted for a new use.
\item[Predation] Eating all or part of another organism (herbivory, carnivory, parasitoidism, parasitism, cannibalism).
\item[Predator satiation] Prey appearing in numbers far beyond predators' capacity, lowering the fraction taken.
\item[Priority effect] The first established competitor excludes the other (unstable L--V case).
\item[Proximate / ultimate cause] How a trait works (mechanism) / why it exists (evolutionary benefit).
\item[$R^*$] Resource level at which a species' growth balances losses; lowest $R^*$ wins.
\item[Rabinowitz rarity] Range $\times$ habitat tolerance $\times$ local population: one common, seven rare types.
\item[Regime shift] Abrupt change to a contrasting state when a threshold is crossed.
\item[Relaxation time] Time scale of return to equilibrium after a disturbance ($1/r$ for the logistic).
\item[Reproductive value ($v$)] Expected future contribution of an individual of a given stage (left eigenvector).
\item[Resilience] Size of the disturbance a state can absorb without shifting (depth of its basin).
\item[Saltationism] Evolution by large sudden jumps rather than gradual change.
\item[Sensitivity] Absolute change in $\lambda$ per absolute change in a matrix entry, $v_iw_j/\mathbf v\cdot\mathbf w$.
\item[SLOSS] Single large or several small reserves; depends on nestedness of species lists.
\item[Stable age (stage) distribution] Constant class proportions reached under fixed vital rates; dominant right eigenvector.
\item[Structural instability] Qualitative behaviour changes under any small change in the equations (basic L--V).
\item[Survivorship curve] $\log l_x$ against age; types I (late loss), II (constant proportion), III (early loss).
\item[Trophic continuum] Modern view of lake trophic states as ranges of $P_{tot}$ at spring circulation.
\item[Turtle excluder device (TED)] Escape hatch in trawl nets protecting large juvenile and older turtles.
\item[Volterra's principle] Killing predator and prey at the same rate raises mean prey and lowers mean predators.
\end{description}
\end{multicols}
}
